\begin{proofbox}
设基准椭圆 $x^2/a^2+y^2/b^2=1$，$a>b>0$，$c^2=a^2-b^2$。
使用\textbf{有符号分母} $D_x=a^2+\lambda$、$D_y=b^2+\lambda$，则
$D_x-D_y=c^2$ 恒定。不要将可能为负的 $D_y$ 写成正半轴长度的平方。
\begin{itemize}
\item $\lambda>-b^2$：两分母为正，半轴为 $\sqrt{D_x},\sqrt{D_y}$，
半焦距平方为 $D_x-D_y=c^2$。
\item $-a^2<\lambda<-b^2$：$D_x>0,D_y<0$，
方程为 $x^2/D_x-y^2/(-D_y)=1$，实、虚半轴平方分别为 $D_x,-D_y$，
半焦距平方为 $D_x+(-D_y)=c^2$。
\end{itemize}
反之，任一同焦点椭圆的长短半轴平方 $A^2,B^2$ 满足 $A^2-B^2=c^2$，
令 $\lambda=A^2-a^2$ 即有 $B^2=b^2+\lambda$。
任一同焦点双曲线满足 $A^2+B^2=c^2$，令同一个 $\lambda=A^2-a^2$，
则 $-B^2=b^2+\lambda$。因此上述两种范围既充分又必要。
\[
\boxed{\frac{x^2}{a^2+\lambda}+\frac{y^2}{b^2+\lambda}=1},\qquad
\lambda>-a^2,\quad\lambda\ne-b^2.
\]
\end{proofbox}
